\textbf{T3: Fabry–Pérot interferometer (10 pts)}

A Fabry–Pérot interferometer consists of two identical parallel planar mirrors
separated by a distance $L$. The space between and outside the mirrors is
filled with air. The mirrors are partially reflective; when light is aimed
towards one of these mirrors along the normal direction, the reflected beam has
intensity $R < 1$ times the intensity of the incident beam. Assume that the
mirrors are symmetric, meaning they interact the same way with light incident
from either side, and lossless. Assume also that they are highly reflective,
meaning $1 - R \ll 1$. A monochromatic laser beam of power $P$ is aimed towards
the interferometer perpendicular to the mirrors. The distance $L$ is chosen so
that the back-reflected beam vanishes, i.e.\ all the optical power is
transmitted through the interferometer.

\begin{center}\includegraphics[width=0.54\linewidth]{figures/EuPhO_2024_Q3_fig1.png}\end{center}

a) (3pts) Show that the laser beam must acquire a nonzero phase shift $\phi$
when it passes through either of the mirrors.

b) (2pts) Find the magnitude of $\phi$.

c) (4pts) At a certain moment, the incident laser beam is switched off rapidly.
Find the total energy of the light that travels back from the interferometer
towards the laser after the laser is switched off.

d) (1pt) Estimate the duration of the light pulse that travels back towards the
laser.
